% Quarantined P03.1 product-obstruction candidate, 30 August 2026.
% Not production verified and not included by the live manuscript.

\begin{proof}[Proof of
\cref{prop:p031-nodal-product-envelope-obstruction-candidate}]
The projection and section satisfy
\[
                         \pi\circ j=\operatorname{id}_U.
\]
Analytic pullback preserves locally Fubini--Study functions, real PL
functions, and semipositivity.  Thus, if \(h\) is admissible for
\(P_U^{\mathrm{raw}}f\), then \(h\circ\pi\) is admissible for \(G\), and
\[
                         G\circ j\geq P_U^{\mathrm{raw}}f.
\]
Conversely, if \(g\) is admissible for \(G\), then \(g\circ j\) is admissible
for \(P_U^{\mathrm{raw}}f\), because \(F\circ j=f\).  Taking suprema gives
\begin{equation}
\label{eq:p031-product-section-identity-candidate}
                 G\circ j=P_U^{\mathrm{raw}}f.
\end{equation}

By \cref{prop:ch15-nodal-local-envelope},
\begin{equation}
\label{eq:p031-product-base-values-candidate}
 P_U^{\mathrm{raw}}f=0\quad\text{on }A\setminus\{o\},
 \qquad
 P_U^{\mathrm{raw}}f=-a\quad\text{on }B.
\end{equation}
Every competitor is bounded above by \(F\), so \(G\leq F\).  Since \(F\) is
continuous, upper-semicontinuous regularization preserves this inequality:
for every \(q\in X_d\),
\[
 E(q)=\limsup_{q'\to q}G(q')
     \leq\limsup_{q'\to q}F(q')
     =F(q).
\]
In particular,
\begin{equation}
\label{eq:p031-product-obstacle-upper-candidate}
                              E\leq F.
\end{equation}

For \(m\geq1\), let \(z_m\in A\setminus\{o\}\) be the rigid point
\[
                       x(z_m)=\varpi^m,\qquad y(z_m)=0.
\]
Then \(z_m\to o\), and
\eqref{eq:p031-product-section-identity-candidate}--%
\eqref{eq:p031-product-base-values-candidate} give
\[
                          G(j(z_m))=0.
\]
It follows from the definition of \(E=G^*\) that \(E(j(o))\geq0\).
On the other hand, \eqref{eq:p031-product-obstacle-upper-candidate} and
\(F(j(o))=f(o)=0\) give \(E(j(o))\leq0\).  Hence
\begin{equation}
\label{eq:p031-product-node-value-candidate}
                              E(j(o))=0.
\end{equation}

The obstacle has \(f(\eta_r)=-a\).  Therefore
\eqref{eq:p031-product-obstacle-upper-candidate} yields
\begin{equation}
\label{eq:p031-product-gauss-upper-candidate}
                         E(j(\eta_r))\leq-a.
\end{equation}

Suppose that \(E\) were classically psh on \(X_d\).  This property is
preserved by analytic pullback and restriction to analytic domains.
Consequently, \(E\circ j\) would be classically psh on the strict closed
disc \(D_r\).  This disc has the singleton Shilov boundary
\(\{\eta_r\}\), so the affinoid maximum principle would imply
\[
                         E(j(o))\leq E(j(\eta_r)).
\]
This contradicts \eqref{eq:p031-product-node-value-candidate} and
\eqref{eq:p031-product-gauss-upper-candidate}, since \(a>0\).  Thus \(E\)
is not classically psh.  Every CLD-v2 psh function is classically psh, so
\(E\) is not CLD-v2 psh either.

Finally,
\[
 X_d=\mathcal M\bigl(
 K\langle x,y,t_1,\ldots,t_d\rangle/(xy)\bigr).
\]
The ideals \((x)\) and \((y)\) are prime, their intersection is \((xy)\),
and their quotients are Tate algebras in \(d+1\) variables.  Thus \((xy)\)
is radical, its two irreducible components both have dimension \(d+1\),
and \(X_d\) is reduced and pure-dimensional.  It is strict by construction.
\end{proof}
